\documentclass[10pt,a4paper]{amsart}
\usepackage[margin=19mm]{geometry}
\usepackage{amsmath,amssymb,mathtools}
\usepackage[scr=rsfs,cal=euler]{mathalfa}
\usepackage{xfrac}
\usepackage[unicode,hidelinks,bookmarks=false]{hyperref}
\newcommand{\ie}{i.e.\ }
\newcommand{\cf}{cf.\ }
\newcommand{\resp}{respectively\ }
\newcommand{\Subsection}{\subsection}
\providecommand{\nobreakdash}{\nobreak\hbox{-}}
\setlength{\emergencystretch}{2em}
\multlinegap0pt
\usepackage[all]{xy}
\usepackage{amssymb}
\usepackage{bm}
\usepackage{xcolor}
\newtheorem{proposition}{Proposition}
\newtheorem{lemma}{Lemma}
\newcommand{\F}{\mathbb F}
\DeclareMathOperator{\Fil}{Fil}
\DeclareMathOperator{\Hom}{Hom}
\newcommand{\Z}{\mathbb Z}
\newcommand{\co}{\mathcal O}
\DeclareMathOperator{\inv}{inv}
\newcommand{\iso}{\cong}
\newcommand{\map}[1]{\xrightarrow{#1}}
\newcommand{\red}{\mathrm{red}}
\title{Rapoport-Zink spaces of type $\mathrm{GU}(2,n-2)$}
\author{Maria Fox, Benjamin Howard,, Naoki Imai}
\date{Source: arXiv:2308.03816v2 (CC BY 4.0)}
\begin{document}
\maketitle

\noindent\emph{Source and licence.} Rapoport-Zink spaces of type $\mathrm{GU}(2,n-2)$. Version arXiv:2308.03816v2 (CC BY 4.0), distributed by arXiv under the Creative Commons Attribution 4.0 International licence, http://creativecommons.org/licenses/by/4.0/, as recorded in the arXiv licence metadata for this exact version and shown on the abstract page https://arxiv.org/abs/2308.03816v2. Full licence notice: \texttt{licenses/arxiv-2308.03816-CC-BY-4.0.txt}.\par\smallskip

\noindent\textit{Editorial scope.} Counted unit: the article's proposition def (source0.tex lines 2280--2299). The article's complete original proof of it is delivered with it (source0.tex lines 2302--2471, one \texttt{proof} environment of 170 lines). No mathematics is written, repaired or re-derived: the statement and the proof are the article's own lines, in the article's own notation, and no retained line is substituted or re-flowed.  Retained uncounted as the source-local context the proof needs: lines 806--811; lines 839--842; lines 875--877; lines 900--902; lines 1110--1123; lines 1202--1207; lines 1215--1217; lines 1326--1351; lines 1389--1394; lines 1413--1415; lines 2209--2213.  Nothing was omitted from any retained span, so the package carries no omission record.  No retained line contains a citation, so the wrapper carries no bibliography and no bibliography entry.  No retained line contains a figure environment, an image-inclusion command or drawing code, so no image is delivered and the package declares no asset dependency.  Editorial formatting only, in addition: an AMS \texttt{amsart} preamble that restates the article's own declarations and macros used by the retained lines, namely the environments \texttt{proposition}, \texttt{lemma} on separate counters, together with the standard packages those lines need.  The retained environments print in this document as Proposition 1, Lemma 1 in order of appearance, whereas the article prints the same items with the numbering of its own full document.  The complete native source of the article as distributed in the arXiv e-print for this exact version is bundled unchanged as \texttt{sources/145328/source0.tex}.
\begin{equation}\label{contraction}
D(\Lambda \otimes \overline{\mathbb{Y}})
\iso 
\breve{\Lambda} \otimes_{\co_{\breve{E}}} D(\overline{\mathbb{Y}}) 
\iso \breve{\Lambda}.
\end{equation}

\begin{equation}\label{Ldef}
 L(X) = D(X)  \stackrel{\varrho_X}{\subset} D(\Lambda \otimes \overline{\mathbb{Y}}) [1/p]
 \stackrel{ \eqref{contraction}} {=} \breve{\Lambda}[1/p] . 
 \end{equation}

\begin{equation}\label{twisted pairing}
b   : \breve{\Lambda}_0  \times \breve{\Lambda}_0 \to \breve{\Z}_p
\end{equation}

\begin{equation}\label{right dual}
L_0^* = \{ x \in  \breve{\Lambda}_0[1/p] :  b(   L_0, x ) \subset \breve{\Z}_p \}
\end{equation}

 \begin{proposition}\label{prop:strata points}
 Under the above bijection,  for any $1\le k \le \lfloor n/2\rfloor$ we have
    \[
\mathrm{RZ}_\Lambda^k( \breve{\F}_p) 
\iso  
\{ L_0  \in \mathrm{RZ}_\Lambda( \breve{\F}_p) : 
\inv(L_0 , \breve{\Lambda}_0) = \lambda_k \} ,
\]
in which 
\[
\lambda_k=( \underbrace{1,\ldots, 1}_{k-1\mathrm{\ times}}, 0,\ldots, 0 , \underbrace{-1,\ldots, -1}_{k\mathrm{\ times}} ) \in \Z^n.
\]
If $k>\lfloor n/2\rfloor$ then $\mathrm{RZ}_\Lambda^k = \emptyset$.
\end{proposition}

\begin{equation}\label{Y tilde functor}
\xymatrix{
{p\Lambda \otimes \overline{\mathbb{Y}}_S  }\ar[r]^a    \ar@/_/[drr]_i  &  { H }  \ar[dr]^b  \ar[r]|{\, d\, }  &  {G}  \ar[r] | {\, d^\vee\, } & {H^\vee}  \ar[r]^{a^\vee}  & {p^{-1}\Lambda \otimes  \overline{\mathbb{Y}}_S  } \\
& &  { \Lambda \otimes  \overline{\mathbb{Y} }_S } \ar[ur]^{b^\vee} \ar@/_/[urr]_{i^\vee}
}
\end{equation}

\begin{equation}\label{Hpol}
\ker(  d^\vee \circ d : H \to H^\vee ) \subset H[p].
\end{equation}

\begin{proposition}\label{prop:RZ tilde points}
There is a   bijection from $Y^k_\Lambda(\breve{\F}_p)$ to the set of pairs of $\breve{\Z}_p$-lattices $(M_0,N_0)$ in $\breve{\Lambda}_0[1/p]$ 
such that
\[
 p\breve{\Lambda}_0 \stackrel{k-1}{\subset}  p M_0^* \stackrel{1}{\subset}  pN_0 \stackrel{n-2k}{\subset}  N_0^* \stackrel{1}{\subset}  M_0 \stackrel{k-1}{\subset} \breve{\Lambda}_0  .
 \]
Here  $M_0^*$ and $N_0^*$ are the right duals, as in \eqref{right dual}, of $M_0$ and $N_0$ with respect to the pairing \eqref{twisted pairing}.
 Under this bijection:
 \begin{enumerate}
 \item
 A point of $Y_\Lambda^k(\breve{\F}_p)$, represented by a diagram \eqref{Y tilde functor}, corresponds to 
 \begin{align*}
 M_0^* &= VD(H^\vee)_1  & N_0^*&= VD(H)_1   \\
   M_0 &= D(H)_0   &     N_0 &= D(H^\vee)_0
 \end{align*}
 all viewed as lattices in $\breve{\Lambda}_0[1/p]$ as in \eqref{Ldef}.
 
 \item
The points of $\mathbf{RZ}^{\le k}_\Lambda(\breve{\F}_p)$ above a given pair $(M_0,N_0)$ are in bijection with $\breve{\Z}_p$-lattices $ L_0 \subset \breve{\Lambda}_0[1/p]$ satisfying
 \[
 L_0^* \stackrel{2}{\subset} L_0\quad \mbox{and}\quad   M_0 \stackrel{k}{\subset} L_0 \stackrel{k-1}{\subset} N_0. 
 \]
 \item
 The points of $\mathbf{RZ}^k_\Lambda(\breve{\F}_p)$ above a given pair $(M_0,N_0)$ are in bijection with $\breve{\Z}_p$-lattices $ L_0 \subset \breve{\Lambda}_0[1/p]$ satisfying the conditions of (2) and, in the notation of  Proposition  \ref{prop:strata points},  $\inv(L_0 , \breve{\Lambda}_0) = \lambda_k$.
 \end{enumerate}
\end{proposition}

\begin{equation}\label{fake minuscule lattices}
\xymatrix{
{p \breve{\Lambda}   }\ar[r]     \ar@/_/[drr]  &  {M} \ar[dr]   \ar[r] & { N_0 \oplus M_1 }  \ar[r] & { N }  \ar[r]  & {p^{-1}  \breve{\Lambda }  } \\
& &  {  \breve{\Lambda} } \ar[ur] \ar@/_/[urr]
}
\end{equation}

\begin{equation}\label{fake dual shift}
\Phi^{-1} N_1 = M_0^* \quad  \mbox{and} \quad  \Phi^{-1} M_1 =N_0^*.
\end{equation}

\begin{align}
\mathscr{V}  &= \mathrm{coker}( \Fil^0 \mathscr{M}_0 \to \Fil^0 \mathscr{N}_0 )  \label{VWdef}  \\
\mathscr{W}  & = \mathrm{coker}(  \mathscr{M}_0 \to  \mathscr{N}_0 ) \nonumber  \\
\mathscr{V}^{(1)} & = \ker(  \mathscr{V} \to \mathscr{W}   ) .  \nonumber
\end{align}

\begin{proposition}\label{prop:beta def}
There is a  unique  $\co_{Y_\Lambda^{k,\red}}$-linear map
\[
\beta : \mathscr{V} \otimes \sigma^*\mathscr{V} \to \co_{Y_\Lambda^{k,\red}}
\]
that, when viewed as a pairing 
\[
\beta : \mathscr{V} \times \mathscr{V} \to \co_{Y_\Lambda^{k,\red}}
\]
(linear in the first variable and $\sigma$-linear in the second),  satisfies the following property:
At a point $s \in Y_\Lambda^{k,\red}(\breve{\F}_p)$,  represented by a pair of lattices $(M_0,N_0)$, 
 as in Proposition \ref{prop:RZ tilde points}, the fiber 
\[
\beta_s : \frac{M_0^*}{N_0^*} \times \frac{M_0^*}{N_0^*} \to \breve{\F}_p
\]
is identified with the reduction  of 
\[
pb : M_0^* \times M_0^* \to  \breve{\Z}_p .
\]
\end{proposition}

\begin{proof}
The uniqueness claim is clear, because we have specified the pairing on fibers.  The existence will follow from the next two lemmas.






\begin{lemma}\label{lem:goof map}
There is an $\co_{Y^\red_\Lambda}$-linear morphism
\[
\sigma^* \Fil^0 \mathscr{N}_0 \to \mathscr{M}_1
\]
whose fiber at any $s\in Y^{k,\red}_\Lambda(\breve{\F}_p)$ is identified with 
\[
 \sigma^* (   M_0^* /  p N_0  )  \map{ 1 \otimes x \mapsto p \Phi x }  M_1 / pM_1  .
\]
Here $M$ and $N$ are the lattices in $\breve{\Lambda}[1/p]$ associated to  $s$ as in \eqref{fake minuscule lattices}.
\end{lemma}

\begin{proof}
Over $Y^{k,\red}_\Lambda$ there is a universal diagram \eqref{Y tilde functor} of $p$-divisible groups.
By \eqref{Hpol}, there is a unique isogeny  $\varpi : H^\vee \to H$  making the diagram 
\[
\xymatrix{
{H^\vee}  \ar[dr]_p  \ar[r]^{\varpi}  & { H  }  \ar[d]^{ d^\vee \circ d}  \\
{    }   &  {   H^\vee  } 
}
\]
commute. This induces an $\co_E$-linear morphism of vector bundles
\[
  \mathscr{N}= \underline{\Hom}_{\co_E}  ( \mathscr{D}(\overline{\mathbb{Y}} ) , \mathscr{D}(H^\vee)  )
  \map{  \varpi  \circ }
   \underline{\Hom}_{\co_E}  ( \mathscr{D}(\overline{\mathbb{Y}} ) , \mathscr{D}(H)  )
   = \mathscr{M},
\]
whose fiber a point of $Y^{k,\red}_\Lambda(\breve{\F}_p)$ is identified with
\[
N / pN \map{p }  M/pM .
\]

Now consider the Frobenius   and Verscheibung morphisms
\[
\mathrm{Fr} :  H \to \sigma^*H
\quad \mbox{and}\quad \mathrm{Ver} : \sigma^* H\to H .
\]
These induce morphisms of vector bundles
\[
 \mathscr{D}( H) \map{ V=\mathscr{D}(  \mathrm{Fr})  }  \sigma^*  \mathscr{D}( H )
 \quad \mbox{and}\quad 
 \sigma^* \mathscr{D}( H ) \map{ F=\mathscr{D}(  \mathrm{Ver})  }   \mathscr{D}( H ),
\]
and similarly with $H$ replaced by $\overline{\mathbb{Y}}$ or $H^\vee$.  
We use this to define a morphism 
\[
\sigma^* \mathscr{N} = 
 \underline{\Hom}_{\co_E} ( \sigma^*\mathscr{D}(\overline{\mathbb{Y}} ) , \sigma^* \mathscr{D}(H^\vee)  ) 
 \map{\alpha}
  \underline{\Hom}_{\co_E}  ( \mathscr{D}(\overline{\mathbb{Y}} )  , \mathscr{D}(H^\vee) )
  = \mathscr{N}
\]
by $x \mapsto  F \circ x \circ V$.
The   fiber of $\alpha$ at a geometric point is identified with (the linearization of)
\[
p\Phi :  N  /pN  \to N / pN.
\]


The maps $\alpha$ and $\varpi \circ $ above sit in a commutative diagram
\[
\xymatrix{
{ \sigma^* \mathscr{N} =  \underline{\Hom}_{\co_E} ( \sigma^*\mathscr{D}(\overline{\mathbb{Y}} ) , \sigma^* \mathscr{D}(H^\vee) ) } \ar[r]^\alpha  \ar[d]_{ F \circ}
&
{  \underline{\Hom}_{\co_E} ( \mathscr{D}(\overline{\mathbb{Y}} ) ,  \mathscr{D}(H^\vee) ) = \mathscr{N} }  \ar[d]^{ \varpi  \circ }   \\
{  \underline{\Hom}_{\co_E} ( \sigma^*\mathscr{D}(\overline{\mathbb{Y}} ) ,   \mathscr{D}(H^\vee) ) } \ar[r]_\beta \ar[d]_{V\circ}
&
{  \underline{\Hom}_{\co_E} ( \mathscr{D}(\overline{\mathbb{Y}} ) ,  \mathscr{D}(H) )   = \mathscr{M} }   \\
{  \sigma^*\mathscr{N}= \underline{\Hom} _{\co_E}( \sigma^*\mathscr{D}(\overline{\mathbb{Y}} ) ,   \sigma^* \mathscr{D}(H^\vee) )  } 
}
\]
in  which the arrow labeled $\beta$ sends $x\mapsto \varpi \circ x \circ V$.
The vertical column on the left is exact, and the image of the final vertical arrow is 
\[
\sigma^* \Fil^0\mathscr{N} \subset \sigma^* \mathscr{N} .
\]
 On fibers, the  composition from the upper left corner to the lower right corner of the square is (the linearization of)
 \[
 N/pN \map{ p^2\Phi}  M/pM .
 \]
 This last map need not be $0$, but the relation $p N_0 \subset N_0^*=\Phi^{-1}M_1$ of Proposition \ref{prop:RZ tilde points} (and \eqref{fake dual shift}, for the equality) implies that its restriction
 \[
 N_0/pN_0 \map{ p^2\Phi}  M_1/pM_1
 \]
vanishes.  Thus the diagram above restricts to a commutative diagram of solid arrows
\[
\xymatrix{
{ \sigma^* \mathscr{N}_0 } \ar[rrr]^\alpha  \ar[d]_{ F \circ}   \ar[drrr]^0&   &  &   {   \mathscr{N}_1 }  \ar[d]^{ \varpi  \circ }   \\
{  \underline{\Hom} ( \sigma^*\mathscr{D}(\overline{\mathbb{Y}} )_0 ,   \mathscr{D}(H^\vee)_1 ) } \ar[rrr]_\beta  \ar[d]_{V\circ}
&  &   &   {  \mathscr{M}_1 }   \\
{  \sigma^* \Fil^0 \mathscr{N}_0 } \ar[d] \ar@{-->}[urrr]  \\
{ 0, }
}
\]
and the exactness of the vertical column on the left implies the existence of a unique dotted arrow making the diagram commute.  This dotted arrow is the morphism we seek.
\end{proof}




\begin{lemma}\label{lem:beta2}
There is a $\co_{Y^{k,\red}_\Lambda}$-linear map
\[
\mathscr{N}_0 \otimes \sigma^* \Fil^0 \mathscr{N}_0 \to \co_{Y^{k,\red}_\Lambda}
\]
whose fiber at any $s\in Y^{k,\red}_\Lambda(\breve{\F}_p)$ is identified with the $\breve{\F}_p$-linear map
\[
\frac{N_0}{pN_0} \otimes \sigma^*\left(    \frac{M_0^*}{pN_0} \right) \to \breve{\F}_p
\]
obtained by reducing $x\otimes (1\otimes y) \mapsto p b(x,y) \in \breve{\Z}_p$ for $x\in N_0$ and $y\in M_0^*$.
\end{lemma}

\begin{proof}
The polarization on  $\overline{\mathbb{Y}}$ induces  a perfect bilinear pairing on the constant vector bundle $\mathcal{D}(\overline{\mathbb{Y}})$, which induces an isomorphism
\[
\mathcal{D}(\overline{\mathbb{Y}})_0 \otimes  \mathcal{D}(\overline{\mathbb{Y}})_1 \iso  \co_{Y^{k,\red}_\Lambda}.
\]
Similarly, there is a  canonical perfect bilinear pairing
\[
\mathcal{D}( H^\vee )_0 \otimes \mathcal{D}(H)_1 \to   \co_{Y^{k,\red}_\Lambda}
\]
on the vector bundles associated to the universal $p$-divisible groups $H$ and $H^\vee$ over $Y^{k,\red}_\Lambda$.  
These induce a perfect bilinear pairing
\begin{align*}
\mathscr{N}_0  \otimes  \mathscr{M}_1
& =  \underline{\Hom}( \mathcal{D}(\overline{\mathbb{Y}})_0 , \mathcal{D}(H^\vee)_0) 
\otimes \underline{\Hom}( \mathcal{D}(\overline{\mathbb{Y}})_1 , \mathcal{D}(H)_1)  \\
& \iso  \underline{\Hom}( \mathcal{D}(\overline{\mathbb{Y}})_0 \otimes \mathcal{D}(\overline{\mathbb{Y}})_1 , \mathcal{D}(H^\vee)_0 \otimes  \mathcal{D}(H)_1)  \\
& \to \co_{Y^{k,\red}_\Lambda}.
\end{align*}
Pulling this  back via the morphism $\sigma^*\Fil^0 \mathscr{N}_0 \to \mathscr{M}_1$ of Lemma \ref{lem:goof map} yields a pairing $\mathscr{N}_0 \otimes \sigma^* \Fil^0 \mathscr{N}_0 \to \co_{Y^{k,\red}_\Lambda}$ having the desired form on  $\breve{\F}_p$-valued points.
\end{proof}


We  now complete the proof of Proposition \ref{prop:beta def}. 
By construction, there is a canonical map $\mathscr{M} \to \mathscr{N}$ respecting filtrations, and we claim that the  composition 
\[
\mathscr{M}_0 \otimes \sigma^* \Fil^0 \mathscr{M}_0 \to \mathscr{N}_0 \otimes \sigma^* \Fil^0 \mathscr{N}_0 \to \co_{Y^{k,\red}_\Lambda}
\]
is trivial. Indeed, taking  fibers  at a point $(M_0,N_0) \in Y_\Lambda^{k,\red}(\breve{\F}_p)$,  the composition is identified with the pairing
\[
\frac{M_0}{pM_0} \times \frac{N_0^*}{pM_0} \to  \frac{N_0}{pN_0} \times \frac{M_0^*}{pN_0}
 \map{pb( - , - )} \breve{\F}_p,
\]
which is trivial as
$
b(M_0,N_0^*) \subset b(M_0,M_0^*) \subset \breve{\Z}_p.
$


Recalling the definitions \eqref{VWdef}, the triviality of the above composition implies that  the morphism of Lemma \ref{lem:beta2}  descends to a morphism
\[
\mathscr{W} \otimes \sigma^* \mathscr{V} \to \co_{Y^{k,\red}_\Lambda}.
\]
The composition
\[
\mathscr{V} \otimes \sigma^*\mathscr{V} \to \mathscr{W} \otimes \sigma^*\mathscr{V} \to \co_{Y_\Lambda^{k,\red}}
\]
(the first arrow is the canonical map $\mathscr{V} \to \mathscr{W}$ on the first tensor factor, and the identity on the second)  
at last defines the desired pairing $\beta$. 
\end{proof}

\end{document}
